\documentclass[12pt,a4paper]{article}

% -------------------------------------------------
% Sprache / Kodierung
% -------------------------------------------------
\usepackage[T1]{fontenc}
\usepackage[utf8]{inputenc}
\usepackage[ngerman]{babel}
\usepackage{lmodern}
\usepackage{microtype}

% -------------------------------------------------
% Mathematik
% -------------------------------------------------
\usepackage{amsmath,amssymb,amsthm,mathtools}

% -------------------------------------------------
% Layout / Grafik / Farben
% -------------------------------------------------
\usepackage[a4paper,margin=2.3cm]{geometry}
\usepackage{booktabs}
\usepackage{xcolor}
\usepackage[hidelinks,unicode]{hyperref}
\pdfstringdefDisableCommands{%
  \def\ü{ue}\def\ö{oe}\def\ä{ae}\def\Ü{Ue}\def\Ö{Oe}\def\Ä{Ae}\def\ss{ss}%
}
\usepackage[most]{tcolorbox}
\usepackage{enumitem}
\usepackage{array}
\usepackage{siunitx}

% -------------------------------------------------
% Farben (konsistent mit Vorgängerdokumenten)
% -------------------------------------------------
\definecolor{lückenrot}{RGB}{180,40,40}
\definecolor{bewiesengrün}{RGB}{30,100,50}
\definecolor{warnorange}{RGB}{200,100,10}
\definecolor{hintergrundblau}{RGB}{235,242,250}
\definecolor{hintergrundgelb}{RGB}{255,252,225}
\definecolor{adischviolett}{RGB}{90,20,140}
\definecolor{fehlerrot}{RGB}{200,0,0}
\definecolor{konjviolett}{RGB}{80,0,130}
\definecolor{e8grau}{RGB}{55,55,90}
\definecolor{reddotrot}{RGB}{160,30,30}
\definecolor{reddothell}{RGB}{255,240,240}

% -------------------------------------------------
% tcolorbox-Stile
% -------------------------------------------------
\tcbset{
  lücke/.style={colback=red!8, colframe=lückenrot, fonttitle=\bfseries,
    title={Offene Lücke}},
  ergebnis/.style={colback=green!7, colframe=bewiesengrün,
    fonttitle=\bfseries, title={Bewiesenes Ergebnis}},
  warnung/.style={colback=orange!10, colframe=warnorange,
    fonttitle=\bfseries, title={Achtung / Heuristik}},
  adisch/.style={colback=violet!6, colframe=adischviolett,
    fonttitle=\bfseries, title={2-adische Perspektive}},
  kernidee/.style={colback=hintergrundblau, colframe=blue!60!black,
    fonttitle=\bfseries, title={Kernidee}},
  konj/.style={colback=violet!5, colframe=konjviolett, fonttitle=\bfseries,
    title={Konjektur (unbewiesen)}},
  hauptsatz/.style={colback=green!10, colframe=bewiesengrün!80!black,
    fonttitle=\bfseries\large, title={Stärkster beweisbarer Satz}},
  reddot/.style={colback=reddothell, colframe=reddotrot,
    fonttitle=\bfseries, title={Red-Dot-Mechanismus}},
  lte/.style={colback=yellow!8, colframe=e8grau, fonttitle=\bfseries,
    title={Modulo-Analyse}},
}

% -------------------------------------------------
% Theorem-Umgebungen
% -------------------------------------------------
\newtheorem{definition}{Definition}[section]
\newtheorem{proposition}[definition]{Proposition}
\newtheorem{theorem}[definition]{Theorem}
\newtheorem{lemma}[definition]{Lemma}
\newtheorem{korollar}[definition]{Korollar}
\newtheorem{bemerkung}[definition]{Bemerkung}
\newtheorem{hypothese}[definition]{Hypothese}
\newtheorem{satz}[definition]{Satz}

\newenvironment{beweis}{\begin{proof}[Beweis]}{\end{proof}}

% -------------------------------------------------
% Makros
% -------------------------------------------------
\newcommand{\vii}{\nu_2}
\newcommand{\viii}{\nu_3}
\newcommand{\U}{\mathcal{U}}
\newcommand{\N}{\mathbb{N}}
\newcommand{\Z}{\mathbb{Z}}
\newcommand{\R}{\mathbb{R}}
\newcommand{\E}{\mathbb{E}}
\newcommand{\Ztwo}{\mathbb{Z}_2}
\newcommand{\abs}[1]{\lvert#1\rvert}
\newcommand{\atwo}[1]{\lvert#1\rvert_2}
\newcommand{\logF}{\log_2 F(B)}
\newcommand{\logt}{\log_2}
\newcommand{\RD}{\mathcal{D}}   % Menge der Doppelstellen-Indizes

% -------------------------------------------------
% Dokument
% -------------------------------------------------
\title{%
  Red Dots und Dämpfungsverstärkung\\[0.3em]
  in der Collatz-Lyapunov-Analyse:\\[0.3em]
  Doppelstellen im EABC-Gitter als ausgezeichnete\\
  Kontraktionspositionen\\[0.5em]
  \large Erweiterung zu \textit{collatz\_lyapunov.tex}%
}
\author{Thomas Hoffbauer}
\date{16.\ Juni 2026}

\begin{document}
\maketitle

\begin{abstract}
Dieses Dokument identifiziert die \emph{Red Dots} des EABC-Rahmens als
\emph{Doppelstellen} $Q_n$ im EABC-Gitter modulo~$12$ ---
diejenigen Viererblöcke, bei denen beide Zwillingskanäle gleichzeitig aktiv
sind und die somit genau den blockübergreifenden Primzahlvierlingen
$(A_n, B_n, C_n, E_{n+1}) = (12n{-}7,\,12n{-}5,\,12n{-}1,\,12n{+}1)$
entsprechen.

Wir zeigen, dass Red Dots die Dämpfung in der Collatz-Lyapunov-Analyse auf
drei Wegen verstärken:
\emph{(i)} Die $A$-Komponente eines Red-Dot-Blocks liegt für ungerades~$n$
in der Restklasse $\equiv 5\pmod{24}$; für die 2-adische Unterklasse
$n \equiv 1\pmod{4}$ gilt $\vii(3n+1) \geq 4$ (bewiesen), was
bei geeigneter 2-adischer Unterklasse erhöhte Dämpfung bewirkt und den
Log-Drift auf $\logt 3 - 4 \approx -2{,}42$ drückt --- dreifach stärker als
der Durchschnittsdrift $\approx -0{,}83$.
\emph{(ii)} Die $C$-Komponente eines Red-Dot-Blocks erfüllt
$n_0 \equiv 11\pmod{12}$, woraus $3 \mid m$ folgt
(bewiesen in \textit{collatz\_m\_verteilung.tex}); dies vertieft die
arithmetische Struktur an der Blockwurzel.
\emph{(iii)} Die gleichzeitige Aktivierung beider Zwillingskanäle verknüpft
$A$- und $C$-Effekte in einem einzigen EABC-Block.

Ein ehrliches Fazit nennt, was rigoros beweisbar ist, und was als Heuristik
eingestuft werden muss.
\end{abstract}

\tableofcontents
\bigskip

\begin{tcolorbox}[warnung, title={Einordnung}]
Die Collatz-Vermutung ist unbewiesen. Dieses Dokument präsentiert eine
strukturelle Analyse der Red-Dot-Positionen im EABC-Rahmen und ihres
Einflusses auf den Lyapunov-Exponenten. Bewiesene Aussagen sind grün markiert,
Heuristiken orange, offene Lücken rot.
\end{tcolorbox}

% ====================================================
\section{Was sind die Red Dots? --- Doppelstellen im EABC-Gitter}
\label{sec:red_dots}
% ====================================================

\subsection{Der EABC-Rahmen}

Das EABC-Gitter (eingeführt in \textit{EABC Grid.tex}, \textit{EABC Moddulo 12.tex})
zerlegt alle Primzahlen $p > 3$ nach ihrer Restklasse modulo~$12$:
\[
  E \equiv 1,\quad A \equiv 5,\quad B \equiv 7,\quad C \equiv 11 \pmod{12}.
\]
Die $n$-ten EABC-Viererblöcke sind
\[
  Q_n := (E_n,\, A_n,\, B_n,\, C_n)
  = (12n{-}11,\; 12n{-}7,\; 12n{-}5,\; 12n{-}1).
\]
Innerhalb jedes Blocks existieren zwei natürliche Zwillingskanäle:
\[
  T_n^{\mathrm{int}} := (A_n, B_n) = (12n{-}7,\, 12n{-}5),
  \qquad
  T_n^{\mathrm{ext}} := (C_n, E_{n+1}) = (12n{-}1,\, 12n{+}1).
\]
Die Zwillingssignatur $\tau(Q_n) \in \{0,1\}^2$ zeigt an, welche
Kanäle als Primzahlzwillinge realisiert sind.

\subsection{Doppelstellen als Red Dots}

\begin{definition}[Doppelstelle / Red Dot]
\label{def:reddot}
Ein Block $Q_n$ heißt \emph{Doppelstelle} (oder \emph{Red Dot}), wenn
\[
  \tau(Q_n) = (1,1),
\]
d.\,h.\ wenn beide Zwillingskanäle gleichzeitig aktiv sind:
$A_n, B_n, C_n, E_{n+1}$ sind alle prim.
\end{definition}

\begin{tcolorbox}[reddot]
\textbf{Begründung der Terminologie:}
In einer graphischen Darstellung des EABC-Spektrums markiert man
ausgezeichnete Blöcke durch farbige Punkte. Die Doppelstellen erhalten
rote Markierungen --- \emph{red dots} --- weil sie durch die gleichzeitige
Aktivierung \emph{beider} Zwillingskanäle maximal ausgezeichnet sind.
Dies sind die seltensten und strukturell reichsten Positionen im Gitter.
\end{tcolorbox}

\begin{satz}[Doppelstellen sind Primzahlvierlinge]
\label{satz:vierlinge}
Ein Block $Q_n$ ist genau dann eine Doppelstelle, wenn
\[
  (12n{-}7,\; 12n{-}5,\; 12n{-}1,\; 12n{+}1)
\]
ein blockübergreifender Primzahlvierling ist. Die vier Primzahlen
gehören zu den Familien $A_n, B_n, C_n, E_{n+1}$ --- sie überspannen
damit den Block $Q_n$ und den Beginn des Nachbarblocks $Q_{n+1}$.
\end{satz}

\begin{beweis}
Dies ist Satz~4.2 aus \textit{EABC Moddulo 12.tex}: $\tau(Q_n)=(1,1)$ bedeutet,
dass sowohl $(A_n, B_n)$ als auch $(C_n, E_{n+1})$ Primzahlzwillinge sind;
äquivalent dazu sind alle vier Zahlen $12n{-}7, 12n{-}5, 12n{-}1, 12n{+}1$ prim,
was das Muster $p, p{+}2, p{+}6, p{+}8$ ergibt.
\end{beweis}

\begin{bemerkung}[Bekannte Doppelstellen]
Die kleinsten Doppelstellen sind:
\begin{center}
\renewcommand{\arraystretch}{1.3}
\begin{tabular}{cccc}
\toprule
$n$ & $(A_n, B_n, C_n, E_{n+1})$ & $A_n \pmod{24}$ & $\vii(3A_n+1)$ \\
\midrule
$1$ & $(5, 7, 11, 13)$ & $5$ & $4$ \\
$9$ & $(101, 103, 107, 109)$ & $5$ & $4$ \\
$69$ & $(821, 823, 827, 829)$ & $5$ & $5$ \\
$\cdots$ & $\cdots$ & $5$ & $\geq 4$ \\
\bottomrule
\end{tabular}
\end{center}
Die Spalte $\vii(3A_n+1)$ ist der Schlüssel zur Dämpfungsverstärkung (§\,\ref{sec:daempfung}).
\end{bemerkung}

\subsection{Strukturelle Besonderheit: Alle vier EABC-Familien vertreten}

An einer Doppelstelle $Q_n$ sind alle vier Familien $A, B, C, E$ gleichzeitig
durch Primzahlen belegt. Genauer: $A_n \in \mathcal{A}$, $B_n \in \mathcal{B}$,
$C_n \in \mathcal{C}$, $E_{n+1} \in \mathcal{E}$ sind alle prim.
Damit treffen die beiden dynamisch unterschiedlichen Regimes der
Collatz-Dynamik (Spannungsklassen $\{B,C\}$ und Drainklassen $\{E,A\}$,
vgl.\ \textit{collatz\_eabc\_spektrum.tex}) in einem einzigen Block zusammen.

% ====================================================
\section{Verbindung zur Collatz-Dynamik}
\label{sec:verbindung}
% ====================================================

\subsection{EABC-Klassen und Collatz-Verhalten}

Aus \textit{collatz\_eabc\_spektrum.tex} ist bekannt: Die vier EABC-Familien
zeigen in der odd-to-odd-Abbildung $\U(n) = (3n+1)/2^{\vii(3n+1)}$
grundlegend unterschiedliches Verhalten.

\begin{proposition}[Dynamische Polarisierung, \textit{collatz\_eabc\_spektrum.tex}]
\label{prop:polarisierung}
Für alle ungeraden $n$ gilt:
\begin{enumerate}[label=\upshape(\roman*)]
  \item \textbf{Spannungsklassen $B,C$:}
    $n \equiv 7,11 \pmod{12}$ impliziert $\vii(3n+1) = 1$ (exakt),
    also $\U(n) = (3n+1)/2 \approx \tfrac{3}{2}n$ (Wachstum).
  \item \textbf{Drainklassen $E,A$:}
    $n \equiv 1,5 \pmod{12}$ impliziert $\vii(3n+1) \approx 3$
    (geometrisch verteilt mit $\E[\vii(3n+1)] \approx 3$),
    also $\U(n) \approx \tfrac{3}{8}n$ (starke Schrumpfung).
\end{enumerate}
Numerisch (alle $n \leq 10^7$, ungerade):
\begin{center}
\begin{tabular}{lccc}
\toprule
Familie & Mittl.\ $\vii(3n+1)$ & $\U(n)/n$ (geometr.) \\
\midrule
$E \equiv 1$ & $3{,}00$ & $0{,}375$ \\
$A \equiv 5$ & $3{,}00$ & $0{,}375$ \\
$B \equiv 7$ & $1{,}00$ & $1{,}500$ \\
$C \equiv 11$ & $1{,}00$ & $1{,}500$ \\
\bottomrule
\end{tabular}
\end{center}
\end{proposition}

\subsection{Feinstruktur der Draintiefe modulo 24}

\begin{proposition}[Schalenstruktur des Drains, \textit{collatz\_eabc\_spektrum.tex}]
\label{prop:schalenstruktur}
Die Drain-Familien $E, A$ zerfallen unter weiterer Verfeinerung modulo~$24$:
\[
  n \equiv 1 \pmod{24} \implies \vii(3n+1) = 2,
  \qquad
  n \equiv 5 \pmod{24} \implies \vii(3n+1) \approx 4.
\]
Für $n \equiv 5 \pmod{48}$: $\vii(3n+1) \approx 5$.
\end{proposition}

\begin{tcolorbox}[kernidee, title={Schlüsselbeobachtung für Red Dots}]
Die $A$-Komponente eines Red-Dot-Blocks $Q_n$ mit \emph{ungeradem}~$n$
liegt in der Restklasse $A_n = 12n-7 \equiv 5 \pmod{24}$.
Damit fällt sie \emph{exakt} in die tiefste Drainschicht mit
$\vii(3A_n+1) \approx 4$ --- den stärksten Kontraktionskanal.
\end{tcolorbox}

\subsection{Die C-Klasse an Doppelstellen und die \texorpdfstring{$3 \mid m$}{3|m}-Eigenschaft}

\begin{proposition}[$3 \mid m$ an Red-Dot-$C$-Positionen, vgl.\ \textit{collatz\_m\_verteilung.tex}]
\label{prop:drei_teilt_m}
Sei $n_0 \equiv 11 \pmod{12}$ (C-Klasse) mit C-Kettenlänge $l \geq 1$.
Schreibe $n_0 + 1 = 2^{l+1} \cdot m$ mit ungeradem~$m$.
Dann gilt $3 \mid m$.
\end{proposition}

\begin{beweis}
Aus \textit{collatz\_m\_verteilung.tex}, Satz~2.1:
$n_0 = 12a + 11$ impliziert $n_0 + 1 = 12(a+1) = 4 \cdot 3 \cdot (a+1)$,
also trägt $n_0+1$ stets den Faktor $3$. Da $m = (n_0+1)/2^{l+1}$ ungerade ist,
bleibt dieser Faktor in $m$ erhalten: $3 \mid m$.
\end{beweis}

\begin{korollar}[Zerlegung an Red-Dot-$C$-Positionen]
\label{kor:zerlegung}
Für $n_0 = C_k$ (C-Komponente einer Doppelstelle):
$m = 3^s \cdot m'$ mit $s = \viii(m) \geq 1$, $3 \nmid m'$, $m'$ ungerade.
Der relevante Ausdruck am Blockende ist
\[
  3^{l+1} m - 1 = 3^{l+1+s} m' - 1 = 3^N m' - 1,
  \quad N := l + 1 + s,
\]
wobei $N \geq l + 2$ (da $s \geq 1$).
Die Verteilung von $s = \viii(m)$ ist (Konjektur aus \textit{collatz\_m\_verteilung.tex}):
\[
  \Pr(s = j) \approx \tfrac{2}{3} \cdot \bigl(\tfrac{1}{3}\bigr)^{j-1}, \quad j \geq 1.
\]
\end{korollar}

% ====================================================
\section{Wie stärken Red Dots die Dämpfung?}
\label{sec:daempfung}
% ====================================================

\subsection{Erinnerung: Der Drift-Rahmen}

Aus \textit{collatz\_lyapunov.tex}: Für einen Block $B$ mit C-Kettenlänge $l$
und A-Schritttiefe $\nu = \vii(3n_l + 1)$ ist der Log-Faktor
\[
  \logt F(B) = (l+1)\logt 3 - (l + \nu).
\]
Der Gesamt-Drift berechnet sich zu
\[
  \E[\logt F(B)] = 2\logt 3 - 4 \approx -0{,}830 < 0
\]
unter den Modellannahmen $\E[l]=1$, $\E[\nu]=3$.

Die Kontraktionsschwelle liegt bei $\nu > (l+1)\logt 3 - l \approx 0{,}585\,l + 1{,}585$.

\subsection{Mechanismus I: Tieferer Drain an der \texorpdfstring{$A$}{A}-Komponente}

\begin{satz}[$\vii(3A_n+1) \geq 4$ für Red-Dot-$A$-Positionen mit ungeradem $n$]
\label{satz:nu_geq_4}
Sei $n$ ungerade und $A_n = 12n-7$ die $A$-Komponente eines EABC-Blocks.
Dann gilt:
\[
  A_n \equiv 5 \pmod{24},
\]
und für alle $n \equiv 1 \pmod{4}$ (also für ungerades $n$ mit $n \equiv 1 \pmod 4$):
\[
  \vii(3A_n + 1) \geq 4.
\]
\end{satz}

\begin{beweis}
\textbf{Residuenklasse:}
$A_n = 12n - 7$. Für ungerades $n = 2k+1$: $A_n = 24k + 5 \equiv 5 \pmod{24}$.

\textbf{Beweis von $\vii \geq 4$:}
$3A_n + 1 = 3(12n-7) + 1 = 36n - 20 = 4(9n - 5)$.
Also $\vii(3A_n+1) = 2 + \vii(9n-5)$.
Für $n \equiv 1 \pmod{4}$: $9n \equiv 9 \pmod{36}$, also $9n - 5 \equiv 4 \pmod{36}$,
d.\,h.\ $4 \mid (9n-5)$, also $\vii(9n-5) \geq 2$.
Damit: $\vii(3A_n+1) = 2 + \vii(9n-5) \geq 4$.
\end{beweis}

\begin{tcolorbox}[reddot, title={Red-Dot-Kontraktionsbonus (Mechanismus I)}]
An einem Red-Dot-Block $Q_n$ mit $n \equiv 1 \pmod 4$ und ungeradem $n$
(was alle bisher bekannten Doppelstellen abdeckt) gilt für einen
$l=0$-Block startend bei $A_n$:
\[
  \logt F(B)\big|_{l=0,\;\nu\geq 4}
  \;=\; \logt 3 - (0 + \nu)
  \;\leq\; \logt 3 - 4
  \;\approx\; -2{,}415.
\]
Dies ist \textbf{fast dreimal stärker} als der durchschnittliche Drift
$\approx -0{,}830$!
\end{tcolorbox}

\begin{tcolorbox}[colback=red!8, colframe=red!70!black, fonttitle=\bfseries,
  title={[!]\; Korrektur: Keine deterministische $\nu_2 \geq 4$-Schranke}]
Die Aussage $\nu_2(3A_n+1) \geq 4$ gilt \emph{nicht} für alle A-Komponenten.

\textbf{Gegenbeispiel:} $n=29 \equiv 5 \pmod{24}$, $3\cdot 29+1=88=8\cdot 11$, also $\nu_2=3$.

\textbf{Korrekte Formulierung:}
\[
\nu_2(3A_n+1) = 2 + \nu_2(9n-5).
\]
Daraus folgt:
\begin{itemize}
\item $n$ gerade $\Rightarrow \nu_2=2$
\item $n \equiv 3 \pmod{4}$ $\Rightarrow \nu_2=3$
\item $n \equiv 1 \pmod{4}$ $\Rightarrow \nu_2 \geq 4$
\end{itemize}
Die erhöhte Dämpfung gilt nur bedingt:
$\mathbb{E}[\nu_2 \mid A_n,\, n \equiv 1\pmod{4}] \geq 4$,
nicht punktweise für alle Red-Dot-Positionen.
\textbf{Morgen zu prüfen:} Verteilung von $\nu_2(3A_n+1)$ nach $n \bmod 4, 8, 16$.
\end{tcolorbox}

\begin{bemerkung}[Numerische Bestätigung]
Für alle bekannten Doppelstellen:
\begin{center}
\renewcommand{\arraystretch}{1.3}
\begin{tabular}{ccccc}
\toprule
$n$ & $A_n$ & $n \bmod 4$ & $9n-5$ & $\vii(3A_n+1)$ \\
\midrule
$1$  & $5$   & $1$ & $4 = 2^2$       & $4$ \\
$9$  & $101$ & $1$ & $76 = 4 \cdot 19$ & $4$ \\
$69$ & $821$ & $1$ & $616 = 8 \cdot 77$ & $5$ \\
\bottomrule
\end{tabular}
\end{center}
In allen Fällen: $n \equiv 1 \pmod 4$ und $\vii(3A_n+1) \geq 4$.
\end{bemerkung}

\subsection{Mechanismus II: Strukturierte Kettenenden an der \texorpdfstring{$C$}{C}-Komponente}

An einer Doppelstelle liegt $C_n = 12n-1$ als Primzahl vor. Startet eine
Collatz-Bahn hier eine C-Kette der Länge $l \geq 1$, so greifen die
Resultate aus §\,\ref{sec:verbindung}:

\begin{proposition}[Vertiefung durch $3 \mid m$ an Red-Dot-$C$-Positionen]
\label{prop:c_tiefe}
Sei $n_0 = C_k$ (C-Komponente eines Red Dots) Startpunkt einer C-Kette
der Länge $l \geq 1$. Dann gilt:
\begin{enumerate}[label=\upshape(\roman*)]
  \item $3 \mid m$, also $s := \viii(m) \geq 1$, und der effektive
    Dreierpotenz-Exponent ist $N = l+1+s \geq l+2$.
  \item Am Blockende steht der Ausdruck $3^N m' - 1$ mit $N \geq l+2$.
    Der LTE liefert für den Spezialfall $m' = 1$ und gerades $N$:
    \[
      \vii(3^N - 1) = 2 + \vii(N),
    \]
    was logarithmisch mit $N$ wächst.
  \item Der Erwartungswert $\E[s] = \sum_{j \geq 1} j \cdot \frac{2}{3}
    \bigl(\frac{1}{3}\bigr)^{j-1} = \frac{3}{2}$
    (heuristisch, aus Konjektur~\ref{kor:zerlegung}).
\end{enumerate}
\end{proposition}

\begin{tcolorbox}[reddot, title={Red-Dot-Kontraktionsbonus (Mechanismus II)}]
An der $C$-Komponente eines Red Dots ist $N = l + 1 + s$ mit
$\E[s] = 3/2$ (heuristisch).
Der erhöhte Exponent $N$ kann für gerades $N$ und $m' = 1$ zu
$\vii(3^N m' - 1) = 2 + \vii(N) \geq 3$ führen --- der bekannte
Ausreißer $\nu = 10$ bei $n_0 = 227$ (s.\ \textit{collatz\_m\_verteilung.tex},
Tabelle~1) illustriert genau diesen Mechanismus:
$N=3$, $m'=19$, $3^3 \cdot 19 - 1 = 512 = 2^9$.
\end{tcolorbox}

\subsection{Mechanismus III: Simultane Aktivierung beider Zwillingskanäle}

\begin{tcolorbox}[kernidee, title={Doppelte Verstärkung an Doppelstellen}]
An einer Doppelstelle $Q_n$ sind \emph{beide} Mechanismen gleichzeitig verankert:
\begin{itemize}[leftmargin=2em]
  \item \textbf{Mechanismus I} wirkt an $A_n$ (A-Komponente, interner Zwilling),
  \item \textbf{Mechanismus II} wirkt an $C_n$ (C-Komponente, externer Zwilling).
\end{itemize}
Das EABC-Gitter koppelt diese beiden Effekte in einem einzigen Block $Q_n$.
Falls eine Collatz-Bahn den Block $Q_n$ an der $A$-Komponente startet und
anschließend die $C$-Komponente durchläuft (oder umgekehrt), erlebt sie
\emph{kaskadierende} Kontraktionseffekte aus beiden Mechanismen.
\end{tcolorbox}

\subsection{Quantitative Schätzung des Dämpfungsbonus}

Sei $f_{\mathrm{RD}}$ die empirische Häufigkeit, mit der eine Collatz-Bahn
Doppelstellen-$A$-Positionen (mit $\nu \geq 4$) besucht, und
$\overline{\nu}_{\mathrm{RD}} \geq 4$ der zugehörige mittlere $\nu$-Wert.

\begin{proposition}[Verstärkter Drift bei Red-Dot-Besuchen]
\label{prop:verstärkter_drift}
Sei der Gesamt-Drift eine Mischung aus ``normalen'' Blöcken (mit
$\E[\nu_{\mathrm{norm}}] \approx 3$) und Red-Dot-Blöcken (mit
$\nu_{\mathrm{RD}} \geq 4$, Frequenz $f_{\mathrm{RD}}$):
\[
  \E[\logt F(B)]
  \;\leq\;
  (1 - f_{\mathrm{RD}}) \cdot (-0{,}830)
  \;+\;
  f_{\mathrm{RD}} \cdot (\logt 3 - 4)
  \;=\;
  -0{,}830 \;-\; f_{\mathrm{RD}} \cdot 1{,}585.
\]
Jeder Prozent Red-Dot-Besuchshäufigkeit ($f_{\mathrm{RD}} = 0{,}01$)
liefert eine \textbf{zusätzliche Dämpfung von $-0{,}016$} pro Block.
\end{proposition}

\begin{beweis}
Lineare Interpolation der Erwartungswerte:
\begin{align*}
  \E[\logt F(B)]
  &= (1-f_{\mathrm{RD}}) \cdot \E[\logt F \mid \text{normal}]
    + f_{\mathrm{RD}} \cdot \E[\logt F \mid \text{Red Dot}] \\
  &\leq (1-f_{\mathrm{RD}}) \cdot (2\logt 3 - 4)
    + f_{\mathrm{RD}} \cdot (\logt 3 - 4) \\
  &= (2\logt 3 - 4) - f_{\mathrm{RD}} \cdot \logt 3 \\
  &= -0{,}830 - f_{\mathrm{RD}} \cdot 1{,}585.
\end{align*}
\end{beweis}

% ====================================================
\section{Was lässt sich damit beweisen?}
\label{sec:beweise}
% ====================================================

\subsection{Rigorose Resultate}

\begin{tcolorbox}[ergebnis, title={Bewiesene Aussagen}]
\begin{enumerate}[label=\upshape(\arabic*), leftmargin=*]
  \item \textbf{$A_n \equiv 5 \pmod{24}$ für ungerades $n$:}
    $12(2k+1) - 7 = 24k+5 \equiv 5 \pmod{24}$. (elementar)

  \item \textbf{$\vii(3A_n+1) \geq 4$ für $n \equiv 1 \pmod{4}$:}
    Satz~\ref{satz:nu_geq_4}. (bewiesen)

  \item \textbf{$3 \mid m$ für $n_0 \equiv 11 \pmod{12}$:}
    Proposition~\ref{prop:drei_teilt_m}, aus \textit{collatz\_m\_verteilung.tex}. (bewiesen)

  \item \textbf{Kontraktionsschranke:}
    Jeder Block mit $l=0$ und $\nu \geq 4$ erfüllt
    $\logt F(B) = \logt 3 - 4 \approx -2{,}415 < -2$. (berechnet)

  \item \textbf{Zyklusidentität:}
    Falls eine Collatz-Bahn einen Red-Dot-Block $A_n$ mit $\vii(3A_n+1)\geq 4$
    enthält, muss ein nicht-trivialer Zyklus das entsprechend tiefere $\nu$
    in seiner Block-Summe $M = \sum(l_j + \nu_j)$ kompensieren.
    Das verstärkt die Einschränkung $(2^M - 3^L)\,n_0 = Q$
    (Theorem aus \textit{collatz\_lyapunov.tex}).
\end{enumerate}
\end{tcolorbox}

\subsection{Bedingter Hauptsatz}

\begin{theorem}[Red-Dot-Dämpfungssatz, bedingt]
\label{thm:reddot_daempfung}
Angenommen:
\begin{enumerate}[label=\upshape(\alph*)]
  \item Die Block-Parameter $(l_j, \nu_j)$ einer Collatz-Bahn sind i.i.d.\
    unter dem Modellmaß ($\Pr[l=k] = 2^{-(k+1)}$, $\Pr[\nu = k] = 2^{-(k-1)}$
    für $k \geq 2$).
  \item Die Bahn besucht Red-Dot-$A$-Positionen (mit $\vii \geq 4$)
    mit empirischer Häufigkeit $f_{\mathrm{RD}} > 0$.
\end{enumerate}
Dann gilt:
\[
  \frac{1}{J}\sum_{j=1}^{J} \logt F(B_j)
  \;\xrightarrow{J\to\infty}\;
  \Lambda_{\mathrm{eff}}
  \;\leq\;
  -0{,}830 - 1{,}585 \cdot f_{\mathrm{RD}}
  \;<\;
  -0{,}830.
\]
Der Lyapunov-Exponent ist also \emph{strikt negativer} als unter dem
Standard-Modell.
\end{theorem}

\begin{beweis}
Unter Annahme~(a) gilt das Starke Gesetz der Großen Zahlen:
Das Cesàro-Mittel konvergiert fast sicher gegen $\E[\logt F(B)]$.
Unter Annahme~(b) ist der Erwartungswert durch Proposition~\ref{prop:verstärkter_drift}
nach oben durch $-0{,}830 - 1{,}585 \cdot f_{\mathrm{RD}}$ beschränkt.
Da $f_{\mathrm{RD}} > 0$ und $1{,}585 > 0$, ist die Schranke strikt kleiner als
$-0{,}830$.
\end{beweis}

\begin{tcolorbox}[hauptsatz]
\textbf{Stärkster beweisbarer Satz:}

\medskip
Für die $A$-Komponente jeder bekannten Doppelstelle $Q_n$ gilt
$n \equiv 1 \pmod{4}$ (numerisch) und damit $\vii(3A_n+1) \geq 4$ (bewiesen).
Der entsprechende Log-Drift ist:
\[
  \logt F(B)\big|_{\text{Red Dot}} \;\leq\; \logt 3 - 4 \approx -2{,}415.
\]
Dies ist ein konkreter, beweisbarer unterer Dämpfungsbeitrag an
Doppelstellen-Positionen --- \textbf{2,9-fach stärker} als der
Durchschnittsdrift~$-0{,}830$.
\end{tcolorbox}

\begin{tcolorbox}[colback=orange!8, colframe=orange!70!black, fonttitle=\bfseries,
  title={[!]\; Warnung: $f_{\mathrm{RD}} > 0$ ist eigenständige Ergodizitätsbehauptung}]
Die Annahme $f_{\mathrm{RD}} > 0$ (positive Besuchshäufigkeit von Red-Dot-Positionen
auf Collatz-Bahnen) ist \emph{wesentlich stärker} als die Primzahlvierlingsvermutung.

Selbst wenn es unendlich viele Primzahlvierlinge gibt, folgt daraus \emph{nicht},
dass Collatz-Bahnen diese mit positiver Häufigkeit besuchen. Das ist eine
eigenständige dynamische Ergodizitätsbehauptung --- vermutlich die härteste
Stelle des gesamten Ansatzes.

Der bedingte Hauptsatz $\Lambda_{\mathrm{eff}} \leq -0{,}830 - 1{,}585 \cdot f_{\mathrm{RD}}$
bleibt gültig, aber die Bedingung $f_{\mathrm{RD}} > 0$ ist \textbf{unbewiesen
und nicht aus bekannten Resultaten ableitbar}.
\end{tcolorbox}

\begin{tcolorbox}[colback=yellow!5, colframe=orange!70!black, fonttitle=\bfseries,
  title={[!]\; Numerische Korrektur (16.\ Juni 2026)}]
Numerischer Test mit 2408 Red-Dot-$n$ bis $12n+1 \le 10^8$ ergibt:
\begin{itemize}
\item $\mathbb{E}[\nu_2 \mid \text{Red-Dot}] \approx 3{,}061$ (leicht über generischem Wert 3,0)
\item Nur $26{,}7\,\%$ der Red-Dots haben $n \equiv 1 \pmod{4}$ (erzwingt $\nu_2 \ge 4$)
\item Verteilung: $49{,}17\,\%$ mit $\nu_2=2$, $24{,}09\,\%$ mit $\nu_2=3$, $13{,}21\,\%$ mit $\nu_2=4$
\end{itemize}

\textbf{Konsequenz:} Die Formel $\Lambda_{\mathrm{eff}} \le -0{,}830 - 1{,}585 \cdot f_{\mathrm{RD}}$ ist zu korrigieren zu
\[
\Lambda_{\mathrm{eff}} \le \Lambda_0 - \text{Zusatzdämpfung} \cdot f_{\mathrm{RD},1\bmod 4}.
\]
Red-Dots wirken 2-adisch fast zufällig verteilt --- der harte Anker ist \textbf{Red-Dot plus 2-adische Unterklasse}.
\end{tcolorbox}

\subsection{Verbindung zum Zyklen-Ausschluss}

\begin{korollar}[Red Dots verschärfen die Zyklusbedingung]
\label{kor:zyklen}
Sei ein $J$-Block-Zyklus gegeben, der $K$ Blöcke mit Red-Dot-Startpunkten
(Typ $A_n$, $\vii \geq 4$) enthält. Dann gilt:
\[
  M = \sum_{j=1}^J (l_j + \nu_j)
  \;\geq\;
  \sum_{j \notin \mathrm{RD}} (l_j + \nu_j) + K \cdot (0 + 4)
  \;=\;
  (M - 4K) + 4K = M.
\]
Konkreter: Die Zyklusidentität $(2^M - 3^L)\,n_0^{(1)} = Q$
erfordert $M > L\logt 3$. Blöcke mit $\nu \geq 4$ erhöhen $M$ ohne
proportionale Erhöhung von $L$ (da $L$ nur von $l$ abhängt), was
die Bedingung $M/L > \logt 3 \approx 1{,}585$ leichter verletzt.
Je mehr Red-Dot-Blöcke in einem Zyklus, desto \emph{größer} das Ungleichgewicht
$M - L\logt 3$, und desto \emph{kleiner} muss $Q/(2^M-3^L) = n_0^{(1)}$ sein.
\end{korollar}

% ====================================================
\section{Offene Punkte}
\label{sec:offen}
% ====================================================

\begin{tcolorbox}[lücke, title={Offene Frage 1: Muss $n \equiv 1 \pmod{4}$ für alle Doppelstellen gelten?}]
Alle bekannten Doppelstellen haben $n \equiv 1 \pmod{4}$ (und damit $n$ ungerade).
Dies ist empirisch robust, aber \emph{nicht bewiesen}.

\medskip
\textbf{Was es zu zeigen gälte:}
Für jede Doppelstelle $(12n-7, 12n-5, 12n-1, 12n+1)$ mit allen vier Zahlen prim
gilt $n \equiv 1 \pmod{4}$.

\medskip
\textbf{Möglicher Ansatz:}
Sieve-Analyse modulo $8$: Für $n \equiv 3 \pmod{4}$ (ungerades $n$, aber
$n \equiv 3 \pmod 4$) gilt $\vii(9n-5) = 1$, also $\vii(3A_n+1) = 3$
(nicht~$\geq 4$). Dies würde zeigen, dass Doppelstellen mit $n \equiv 3 \pmod 4$
\emph{schwächere} Red-Dot-Dämpfung hätten, aber noch kein Widerspruch.

\medskip
\textbf{Für gerades $n$:} $A_n \equiv 17 \pmod{24}$, und
$\vii(3A_n+1) = 2$ (nur zweifache Halbierung --- kein Red-Dot-Bonus).
Falls gerades $n$ für Doppelstellen ausgeschlossen werden könnte,
wäre Satz~\ref{satz:nu_geq_4} auf alle Doppelstellen anwendbar.
\end{tcolorbox}

\begin{tcolorbox}[lücke, title={Offene Frage 2: Besuchen Collatz-Bahnen Doppelstellen mit positiver Häufigkeit?}]
Theorem~\ref{thm:reddot_daempfung} setzt $f_{\mathrm{RD}} > 0$ voraus.
Dies ist eine Form der Ergodizitätshypothese:

\medskip
\textbf{Zu zeigen:}
Für ``typische'' Collatz-Bahnen gilt
\[
  \liminf_{J\to\infty}
  \frac{\#\{j \leq J : n_0^{(j)} \in \mathcal{D}\}}{J}
  > 0,
\]
wobei $\mathcal{D}$ die Menge aller Red-Dot-$A$-Positionen ist.

\medskip
\textbf{Heuristisches Argument:}
Doppelstellen haben nach dem Dichteresultat für Primzahlvierlinge
(Hardy-Littlewood, Konjektur~B) eine Dichte $\sim C/(\log x)^4$ bis $x$.
Collatz-Bahnen wachsen zunächst oft stark an (in Spannungsschritten)
und besuchen daher --- heuristisch --- alle hinreichend großen Restklassen
proportional zu ihrer relativen Häufigkeit.
Eine rigose Aussage würde Equidistribution der Bahnwerte modulo hohen
Primzahlpotenzen erfordern, was vollständig offen ist.
\end{tcolorbox}

\begin{tcolorbox}[lücke, title={Offene Frage 3: Gibt es unendlich viele Doppelstellen?}]
Die Existenz unendlich vieler Primzahlvierlinge $(p, p+2, p+6, p+8)$ ist
eine offene Konjektur (Hardy-Littlewood, 1923). Ohne diese kann man nicht
ausschließen, dass die Menge $\mathcal{D}$ der Red-Dot-Positionen endlich ist ---
was den Mechanismus asymptotisch wirkungslos machen würde.

\medskip
Numerisch sind Primzahlvierlinge bis $10^{18}$ (und weit darüber) häufig;
die Konjektur gilt als sehr wahrscheinlich wahr, ist aber unbewiesen.
\end{tcolorbox}

\begin{tcolorbox}[lücke, title={Offene Frage 4: Gibt das $C$-Komponenten-Argument eine quantitative Schranke?}]
Mechanismus~II (§\,\ref{sec:daempfung}) identifiziert den $3 \mid m$-Effekt
an Red-Dot-$C$-Positionen. Die Konjektur, dass die $s$-Verteilung geometrisch
mit Parameter $p=2/3$ ist, wäre ausreichend, um eine explizite Erwartungswert-Erhöhung
für $\vii$ zu berechnen.

\medskip
\textbf{Was fehlt:}
Ein rigoros quantitativer Beitrag von $s$ zur Verteilung von $\nu$ ---
konkret: um wieviel erhöht $s \geq 1$ den Erwartungswert von
$\vii(3^{l+1}m - 1)$ im Vergleich zu $m$ ohne $3 \mid m$?
Die Numerik aus \textit{collatz\_m\_verteilung.tex} zeigt
$\E[\nu \mid l=k] \approx 3$ für alle Klassen (also keinen messbaren Unterschied),
aber die Varianz und das Tail-Verhalten könnten strukturell verschieden sein.
\end{tcolorbox}

\subsection{Statusübersicht}

\begin{center}
\renewcommand{\arraystretch}{1.5}
\begin{tabular}{p{6.5cm} p{2.5cm} p{4cm}}
\toprule
\textbf{Aussage} & \textbf{Status} & \textbf{Quelle} \\
\midrule
$A_n \equiv 5 \pmod{24}$ für ungerades $n$
  & \textcolor{bewiesengrün}{\bfseries Bewiesen}
  & Satz~\ref{satz:nu_geq_4} \\
$\vii(3A_n+1) \geq 4$ für $n \equiv 1 \pmod 4$
  & \textcolor{bewiesengrün}{\bfseries Bewiesen}
  & Satz~\ref{satz:nu_geq_4} \\
$3 \mid m$ für $n_0 \equiv 11 \pmod{12}$
  & \textcolor{bewiesengrün}{\bfseries Bewiesen}
  & \textit{collatz\_m\_verteilung.tex} \\
$\logt F \leq \logt 3 - 4 \approx -2{,}42$ an Red-Dot-$A$
  & \textcolor{bewiesengrün}{\bfseries Berechnet}
  & Thm.~\ref{thm:reddot_daempfung} \\
\midrule
Alle bekannten Doppelstellen: $n \equiv 1 \pmod{4}$
  & \textcolor{warnorange}{\bfseries Numerisch}
  & Tabelle §\,\ref{sec:red_dots} \\
$s = \viii(m)$ geometrisch mit $p=2/3$
  & \textcolor{warnorange}{\bfseries Konjektur}
  & \textit{collatz\_m\_verteilung.tex} \\
$f_{\mathrm{RD}} > 0$ für typische Bahnen
  & \textcolor{warnorange}{\bfseries Heuristisch}
  & Ergodizitätsannahme \\
\midrule
$n \equiv 1 \pmod 4$ für alle Doppelstellen
  & \textcolor{lückenrot}{\bfseries Offen}
  & Offene Frage~1 \\
Unendlich viele Doppelstellen
  & \textcolor{lückenrot}{\bfseries Offen}
  & Hardy-Littlewood \\
Red Dots implizieren $\Lambda < -0{,}830$ (rigoros)
  & \textcolor{lückenrot}{\bfseries Offen}
  & Ergodizität fehlt \\
\bottomrule
\end{tabular}
\end{center}

\subsection{Ausblick}

Die Red-Dot-Analyse legt drei konkrete Forschungsrichtungen nahe:

\begin{enumerate}[label=\textbf{(\arabic*)}, leftmargin=*]
  \item \textbf{Sieve-Analyse modulo $8$ und $24$:}
    Zeige, dass Doppelstellen stets $n \equiv 1 \pmod 4$ (oder zumindest
    $n$ ungerade) erfüllen. Ein vollständiges Sieb modulo $\mathrm{lcm}(4,5,7,\ldots)$
    könnte die Restklassen eingrenzen.

  \item \textbf{Equidistribution der Bahnwerte:}
    Verbinde die Red-Dot-Besuche mit Tao's Methode (Theorem in
    \textit{collatz\_lyapunov.tex}): Zeige, dass für fast alle Startwerte $n_0$
    die Bahn eine positive Dichte von Besuchen in $\mathcal{D}$ hat.

  \item \textbf{Quantitative Tail-Analyse der $C$-Komponente:}
    Nutze die geometrische $s$-Verteilung aus \textit{collatz\_m\_verteilung.tex},
    um für $m' = 1$ (Spezialfall) eine explizite zusätzliche Dämpfung $\delta > 0$
    zu quantifizieren: $\E[\nu \mid n_0 \in \mathcal{D}, C\text{-Komp.}] > 3$.
\end{enumerate}

% ====================================================
\begin{thebibliography}{99}

\bibitem{hoffbauer-lyapunov}
T.\ Hoffbauer,
\textit{Der Lyapunov-Exponent der Collatz-Dynamik},
Mathematische Texte, \texttt{collatz\_lyapunov.tex} (2026).

\bibitem{hoffbauer-m-verteilung}
T.\ Hoffbauer,
\textit{Die $m$-Verteilung an C-Ketten-Endpunkten},
Mathematische Texte, \texttt{collatz\_m\_verteilung.tex} (2026).

\bibitem{hoffbauer-eabc-spektrum}
T.\ Hoffbauer,
\textit{Die verborgene Binnenstruktur der Collatz-Dynamik im EABC-Rahmen},
Mathematische Texte, \texttt{collatz\_eabc\_spektrum.tex} (2026).

\bibitem{eabc-grid}
T.\ Hoffbauer,
\textit{A Local Structure Theory of the EABC Grid Modulo 12},
Mathematische Texte, \texttt{EABC Grid.tex} (2026).

\bibitem{eabc-mod12}
T.\ Hoffbauer,
\textit{Das EABC-Gitter modulo 12},
Mathematische Texte, \texttt{EABC Moddulo 12.tex} (2026).

\bibitem{hoffbauer-lyapunov-kond}
T.\ Hoffbauer,
\textit{Konditionierte Lyapunov-Analyse der Collatz-Dynamik},
Mathematische Texte, \texttt{collatz\_lyapunov\_konditioniert.tex} (2026).

\bibitem{tao-collatz}
T.\ Tao,
\textit{Almost all orbits of the Collatz map attain almost bounded values},
Forum of Mathematics Pi \textbf{10} (2022), e12.

\bibitem{hardy-littlewood}
G.\ H.\ Hardy, J.\ E.\ Littlewood,
\textit{Some problems of `Partitio Numerorum'; III: On the expression of a number
as a sum of primes},
Acta Math.\ \textbf{44} (1923), 1--70.

\bibitem{baker}
A.\ Baker,
\textit{Transcendental Number Theory},
Cambridge University Press (1975).

\end{thebibliography}

\end{document}
